\documentclass{article}
\usepackage{amsmath}
\usepackage{amssymb}
\usepackage{siunitx}

\begin{document}

\section*{Choosing the Out-of-Plane Lattice Vector for a 2D System}

When simulating a truly two-dimensional material using a three-dimensional code with periodic boundary conditions, one must pick a third lattice vector~$\mathbf{c}$ that

\begin{itemize}
  \item is orthogonal to the in-plane vectors $\mathbf{a}$ and $\mathbf{b}$,
  \item has a length large enough to avoid spurious interactions between periodic images,
  \item and ideally preserves a consistent “effective density” so that any algorithms expecting a 3D cell volume still behave reasonably.
\end{itemize}

Here’s one way to pick $\mathbf{c}$:

\subsection*{1.\quad Relate cell volume and particle radius in 3D}

In three dimensions, if we imagine each of the $N$ particles occupies a volume equal to that of a sphere of radius~$r_\mathrm{avg}$, then
\begin{equation}
  V_\mathrm{cell}
  = N \, V_{\rm sphere}
  = N \,\frac{4}{3}\pi r_\mathrm{avg}^3
  \,.
  \label{eq:3D-volume}
\end{equation}

\subsection*{2.\quad Relate cell area and particle radius in 2D}

In two dimensions, each particle “occupies” a circle of radius~$r_\mathrm{avg}$, so
\begin{equation}
  A_\mathrm{cell}
  = N \, A_{\rm circle}
  = N \,\pi r_\mathrm{avg}^2
  \,.
  \label{eq:2D-area}
\end{equation}

Since $A_\mathrm{cell}$ and $N$ are known, solve \eqref{eq:2D-area} for
\begin{equation}
  r_\mathrm{avg}
  = \sqrt{\frac{A_\mathrm{cell}}{N\,\pi}}.
  \label{eq:r-avg}
\end{equation}

\subsection*{3.\quad Compute an effective 3D cell volume}

Plugging \eqref{eq:r-avg} into \eqref{eq:3D-volume}, we obtain an effective 3D volume:
\begin{equation}
  V_\mathrm{cell}
  = N \,\frac{4}{3}\,\pi
    \left(\sqrt{\frac{A_\mathrm{cell}}{N\,\pi}}\right)^{\!3}
  = \frac{4}{3}\,\pi
    \,\frac{A_\mathrm{cell}^{3/2}}{(N\,\pi)^{1/2}}
  \,.
  \label{eq:V-cell-effective}
\end{equation}

\subsection*{4.\quad Determine the out-of-plane lattice constant}

By definition of the parallelepiped volume,
\[
  V_\mathrm{cell}
  = \bigl|\mathbf{a}\times\mathbf{b}\bigr|\;|\mathbf{c}|
  = A_\mathrm{cell}\;c
  \,,
\]
since $\mathbf{c}\perp\mathbf{a},\mathbf{b}$ by construction.  Hence
\begin{equation}
  c \;=\; |\mathbf{c}|
  \;=\;
  \frac{V_\mathrm{cell}}{A_\mathrm{cell}}
  = \frac{\displaystyle
      \frac{4}{3}\,\pi
      \,A_\mathrm{cell}^{3/2}/(N\,\pi)^{1/2}
    }{A_\mathrm{cell}}
  = \frac{4}{3}\,\pi
    \,\frac{\sqrt{A_\mathrm{cell}}}{\sqrt{N\,\pi}}
  = \frac{4}{3}\,\pi\,r_\mathrm{avg}
  \,.
  \label{eq:c-length}
\end{equation}

\subsection*{5.\quad Algorithmic summary}

\begin{enumerate}
  \item Compute the 2D cell area
    \[
      A_\mathrm{cell} \;=\; \bigl|\mathbf{a}\times\mathbf{b}\bigr|.
    \]
  \item Compute
    \(\displaystyle r_\mathrm{avg} = \sqrt{A_\mathrm{cell}/(N\,\pi)}\).
  \item Compute the effective 3D cell volume via
    \(\displaystyle V_\mathrm{cell} = N\,(4/3)\,\pi\,r_\mathrm{avg}^3\).
  \item Set
    \(\displaystyle \mathbf{c} = c\,\hat{\mathbf{n}}\),
    where \(\hat{\mathbf{n}}\) is the unit vector perpendicular to
    \(\mathbf{a}\) and \(\mathbf{b}\), and
    \(\displaystyle c = V_\mathrm{cell}/A_\mathrm{cell}\).
  \item Now run your regular 3D routines using the cell
    \(\{\mathbf{a},\mathbf{b},\mathbf{c}\}\).
\end{enumerate}

\textbf{Note.}  This construction assumes a perfectly flat 2D lattice
(i.e.\ well-defined $A_\mathrm{cell}$) and aims only to give a
self-consistent “height” for $\mathbf{c}$.  In practice you may still
want to increase $c$ further to avoid any long-range image interactions
(e.g.\ electrostatics, van der Waals).

\subsection*{6.\quad Closed-form expression for \(c\)}

Since 
\[
A_{\rm cell} = \lvert \mathbf{a}\times\mathbf{b}\rvert 
             = a\,b\,\cos\theta,
\]
we have from step 4
\[
c 
= \frac{V_{\rm cell}}{A_{\rm cell}}
= \frac{\displaystyle\frac{4}{3}\,\pi\,r_{\rm avg}^3\,N}{A_{\rm cell}}
= \frac{4}{3\sqrt{\pi}}
    \sqrt{\frac{A_{\rm cell}}{N}}
= \frac{4}{3\sqrt{\pi}}
    \sqrt{\frac{a\,b\,\cos\theta}{N}}.
\]
In the common case of an orthogonal basis (\(\theta=90^\circ\), \(\cos\theta=1\)),
\[
\boxed{
c 
= z 
= \frac{4}{3\sqrt{\pi}}
  \,\sqrt{\frac{a\,b}{N}}
}.
\]

\end{document}
